\section*{Capítulo \ref{ch:g12:cells-and-electrolysis} --- Pilas electroquímicas y electrólisis}

\begin{solution}{exo:g12:cells-and-electrolysis:1}
El electrodo de cinc es el \omterm{def:g12:cells-and-electrolysis:anode}{ánodo}, \ce{Zn -> Zn^{2+} + 2e-}; el electrodo
de cobre es el \omterm{def:g12:cells-and-electrolysis:anode}{cátodo}, \ce{Cu^{2+} + 2e- -> Cu}.
\end{solution}

\begin{solution}{exo:g12:cells-and-electrolysis:2}
% ledger: const:F
$Q = 0.20 \times 1800 = \qty{360}{C}$; $n(e^-) = 360/96485 =
\qty{3.7e-3}{mol}$.
\end{solution}

\begin{solution}{exo:g12:cells-and-electrolysis:3}
Cierra el circuito y mantiene eléctricamente neutra cada disolución, ya
que sus \omterm{def:g9:ions:ion}{iones} pasan a los compartimentos. Sin él, el circuito está
abierto: no circula corriente, el voltímetro marca cero y una lámpara
sigue apagada.
\end{solution}

\begin{solution}{exo:g12:cells-and-electrolysis:4}
$2.5 \times 3600 = \qty{9000}{C}$.
\end{solution}

\begin{solution}{exo:g12:cells-and-electrolysis:5}
No: la reacción va en el sentido opuesto al espontáneo. El generador
aporta la energía.
\end{solution}

\begin{solution}{exo:g12:cells-and-electrolysis:6}
\omterm{def:g12:cells-and-electrolysis:anode}{Ánodo} (cobre, $-$): \ce{Cu -> Cu^{2+} + 2e-}. \omterm{def:g12:cells-and-electrolysis:anode}{Cátodo} (plata, $+$):
\ce{Ag+ + e- -> Ag}. Global: \ce{Cu + 2Ag+ -> Cu^{2+} + 2Ag}.
\end{solution}

\begin{solution}{exo:g12:cells-and-electrolysis:7}
Tres años son $3 \times 365 \times 24 = \qty{26280}{h}$;
$0.010 \times 26280 = \qty{263}{mA.h}$.
\end{solution}

\begin{solution}{exo:g12:cells-and-electrolysis:8}
% ledger: const:F, aw:Ag
$Q = 0.50 \times 2400 = \qty{1200}{C}$; $n(e^-) = n(\ce{Ag}) =
1200/96485 = \qty{0.0124}{mol}$; $m = 0.0124 \times 107.9 =
\qty{1.34}{g}$.
\end{solution}

\begin{solution}{exo:g12:cells-and-electrolysis:9}
La corriente va por el cable del cobre ($+$) al cinc ($-$), en sentido
opuesto a los \omterm{def:g9:inside-the-atom:nucleus}{electrones}. Los \omterm{def:g9:ions:ion}{iones} \ce{K+} se desplazan hacia el
compartimento del cobre, que pierde \ce{Cu^{2+}}; los \omterm{def:g9:ions:ion}{iones} \ce{NO3-},
hacia el compartimento del cinc, que gana \ce{Zn^{2+}}.
\end{solution}

\begin{solution}{exo:g12:cells-and-electrolysis:10}
En ambos casos, el \omterm{def:g12:cells-and-electrolysis:anode}{ánodo} es donde ocurre la \omterm{def:g11:redox:oxidation}{oxidación}. En una pila, el
\omterm{def:g12:cells-and-electrolysis:anode}{ánodo} es $-$, la reacción es espontánea y la energía química se
convierte en energía eléctrica; en una \omterm{def:g12:cells-and-electrolysis:electrolysis}{electrólisis}, el \omterm{def:g12:cells-and-electrolysis:anode}{ánodo} está unido
al $+$ del generador, la reacción es forzada y la energía eléctrica se
convierte en energía química.
\end{solution}

\begin{solution}{exo:g12:cells-and-electrolysis:11}
\qty{12}{mL}: la ecuación \ce{2H2O -> 2H2 + O2} da dos \omterm{def:g7:atoms-and-molecules:molecule}{moléculas} de
dihidrógeno por cada una de dioxígeno, y cantidades iguales de gas
ocupan volúmenes iguales.
\end{solution}

\begin{solution}{exo:g12:cells-and-electrolysis:12}
% ledger: aw:Zn, const:F
$n(\ce{Zn}) = 5.0/65.4 = \qty{0.076}{mol}$;
$n(\ce{Cu^{2+}}) = 0.100 \times 0.50 = \qty{0.050}{mol}$; reaccionan
\omterm{def:g10:the-mole:amount}{mol} a \omterm{def:g10:the-mole:amount}{mol}, así que los \omterm{def:g9:ions:ion}{iones} cobre se agotan primero.
$n(e^-) = 2 \times 0.050
= \qty{0.100}{mol}$, $Q = \qty{9.65e3}{C} = \qty{2.68}{A.h}$. A \qty{50}{mA}: $2.68/0.050 = \qty{54}{h}$.
\end{solution}

\begin{solution}{exo:g12:cells-and-electrolysis:13}
$Q = 2.0 \times 3600 = \qty{7200}{C}$; $E = 7200 \times 1.5 =
\qty{1.08e4}{J}$, es decir, $1.5 \times 2.0 = \qty{3.0}{W.h}$.
\end{solution}

\begin{solution}{exo:g12:cells-and-electrolysis:14}
% ledger: const:F
$Q = \qty{7200}{C}$, $n(e^-) = \qty{0.0746}{mol}$. Dos \omterm{def:g9:inside-the-atom:nucleus}{electrones} por
\ce{H2}: $n(\ce{H2}) = \qty{0.0373}{mol}$, $V = \qty{0.90}{L}$. Cuatro por
\ce{O2}: $n(\ce{O2}) = \qty{0.0187}{mol}$, $V = \qty{0.45}{L}$.
\end{solution}

\begin{solution}{exo:g12:cells-and-electrolysis:15}
El cobre disuelto en el \omterm{def:g12:cells-and-electrolysis:anode}{ánodo} es igual al cobre depositado, ya que por
ambos pasan los mismos \omterm{def:g9:inside-the-atom:nucleus}{electrones}: $2.84/3.00 = \qty{94.7}{\percent}$ de
cobre.
\end{solution}

\begin{solution}{pb:g12:cells-and-electrolysis:1}
% ledger: const:F, const:NA, aw:Li, dch:CH4
\textbf{1.} El \omterm{def:g12:cells-and-electrolysis:anode}{ánodo}: allí se oxida el litio.

\textbf{2.} El electrodo de grafito, el \omterm{def:g12:cells-and-electrolysis:anode}{ánodo}.

\textbf{3.} Los \omterm{def:g9:inside-the-atom:nucleus}{electrones} atraviesan el teléfono del electrodo de
grafito al electrodo de óxido; los \omterm{def:g9:ions:ion}{iones} litio cruzan el electrolito en
el mismo sentido, del grafito al óxido.

\textbf{4.} Un cargador puede obligar a su reacción a ir en sentido
inverso.

\textbf{5.} \qty{4.000}{A.h}, es decir, $4.000 \times 3600 =
\qty{14400}{C}$.

\textbf{6.} $14400/96485 = \qty{0.1492}{mol}$.

\textbf{7.} $0.1492 \times \num{6.022e23} = \num{8.99e22}$ \omterm{def:g9:inside-the-atom:nucleus}{electrones}.

\textbf{8.} $4.000/0.25 = \qty{16}{h}$.

\textbf{9.} $0.20 \times 14400 = \qty{2880}{C}$.

\textbf{10.} $E = 14400 \times 3.7 = \qty{5.3e4}{J}$, unos \qty{53}{kJ}.

\textbf{11.} $53280/3600 = \qty{14.8}{W.h}$.

\textbf{12.} Quemar \qty{1}{g} de metano libera algo más de energía
(\qty{55.7}{kJ}) de la que almacena una batería de teléfono llena.

\textbf{13.} $1000/14.8 = 67.6$: 67 cargas completas.

\textbf{14.} \ce{Li+ + e- -> Li}: una \omterm{def:g11:redox:oxidation}{reducción}.

\textbf{15.} El cargador obliga a la reacción a ir en el sentido opuesto
al espontáneo.

\textbf{16.} $14400/2.0 = \qty{7200}{s}$, \qty{2.0}{h}.

\textbf{17.} Un \omterm{def:g9:ions:ion}{ion} litio por \omterm{def:g9:inside-the-atom:nucleus}{electrón}: \qty{0.1492}{mol}.

\textbf{18.} $0.1492 \times 6.9 = \qty{1.03}{g}$.

\textbf{19.} Alrededor de \qty{1.03}{g} de litio va y viene entre los
electrodos en una carga completa.
\end{solution}
